\documentclass[10pt,a4paper]{amsart}
\usepackage[margin=19mm]{geometry}
\usepackage{amsmath,amssymb,mathtools}
\usepackage[scr=rsfs,cal=euler]{mathalfa}
\usepackage{xfrac}
\usepackage[unicode,hidelinks,bookmarks=false]{hyperref}
\newcommand{\ie}{i.e.\ }
\newcommand{\cf}{cf.\ }
\newcommand{\resp}{respectively\ }
\newcommand{\Subsection}{\subsection}
\providecommand{\nobreakdash}{\nobreak\hbox{-}}
\setlength{\emergencystretch}{2em}
\multlinegap0pt
\usepackage{amssymb}
\usepackage[shortlabels]{enumitem}
\usepackage{tikz}
\usepackage{cancel}
\newtheorem{lemma}{Lemma}
\newcommand{\mptn}{\mathcal{S}} %%% pattern manifold
\newcommand{\trans}{^\top}
\title{A strong nullity parameter for rooted graphs}
\author{Aida Abiad, Mary Flagg, H. Tracy Hall, Jephian C.-H. Lin, Bryan Shader}
\date{Source: arXiv:2603.00899v1 (CC BY 4.0)}
\begin{document}
\maketitle

\noindent\emph{Source and licence.} A strong nullity parameter for rooted graphs. Version arXiv:2603.00899v1 (CC BY 4.0), distributed by arXiv under the Creative Commons Attribution 4.0 International licence, http://creativecommons.org/licenses/by/4.0/, as recorded in the arXiv licence metadata for this exact version and shown on the abstract page https://arxiv.org/abs/2603.00899v1. Full licence notice: \texttt{licenses/arxiv-2603.00899-CC-BY-4.0.txt}.\par\smallskip

\noindent\textit{Editorial scope.} Counted unit: the article's lemma lem:schursnip (source0.tex lines 684--687). The article's complete original proof of it is delivered with it (source0.tex lines 688--745, one \texttt{proof} environment of 58 lines). No mathematics is written, repaired or re-derived: the statement and the proof are the article's own lines, in the article's own notation, and no retained line is substituted or re-flowed.  No further source-local context is needed: every cross-reference of every retained line resolves inside the two delivered spans.  Nothing was omitted from any retained span, so the package carries no omission record.  No retained line contains a citation, so the wrapper carries no bibliography and no bibliography entry.  No retained line contains a figure environment, an image-inclusion command or drawing code, so no image is delivered and the package declares no asset dependency.  Editorial formatting only, in addition: an AMS \texttt{amsart} preamble that restates the article's own declarations and macros used by the retained lines, namely the environments \texttt{lemma} on separate counters, together with the standard packages those lines need.  The retained environments print in this document as Lemma 1 in order of appearance, whereas the article prints the same items with the numbering of its own full document.  The complete native source of the article as distributed in the arXiv e-print for this exact version is bundled unchanged as \texttt{sources/195525/source0.tex}.
\begin{lemma}
\label{lem:schursnip}
Let $G$ be a graph and $A\in\mptn(G)$.  Suppose $\alpha\subseteq V(G)$ is a subset of vertices such that $A[\alpha]$ is invertible and $A/A[\alpha]\in\mptn(H)$ for some graph $H$ in which $N(\alpha)$ is a clique.  Then $A$ having SAP implies that $A/A[\alpha]$ has SAP, and for any given $i\notin\alpha$, $A$ having $i$-SNIP implies that $A/A[\alpha]$ has $i$-SNIP.  
\end{lemma}

\begin{proof}
Let $a = |\alpha|$ and $b = |N(\alpha)|$.  Without loss of generality, we may assume the first rows/columns are corresponding to $\alpha$ and write 
\[
    A = 
        \begin{bmatrix}
        Q & B\trans \\
        B & C
    \end{bmatrix}
\]
with $Q = A[\alpha]$.  We may further assume the first $b$ rows in $B$ correspond to $N(\alpha)$, so $B$ is zero except for the first $b$ rows.  By the Schur complement, we have 
\[
    \begin{bmatrix}
        I & O \\
        -BQ^{-1} & I
    \end{bmatrix}
    \begin{bmatrix}
        Q & B\trans \\
        B & C
    \end{bmatrix}
    \begin{bmatrix}
        I & -Q^{-1}B\trans \\
        O & I
    \end{bmatrix} = 
    \begin{bmatrix}
    Q & O \\
    O & A / Q
    \end{bmatrix},
\]
which we write as $E\trans A E = Q\oplus (A/Q)$.

Let $X$ be a real symmetric matrix such that $A/Q\circ X = I\circ X = O$.  Define  
\[
    \widehat{X} = E(O\oplus X)E\trans = 
    \begin{bmatrix}
        Q^{-1}B\trans XBQ^{-1} & -Q^{-1}B\trans X\\
        -XBQ^{-1} & X
    \end{bmatrix}.
\]
We claim that $A\circ\widehat{X} = I\circ\widehat{X} = O$.  
Since $A/Q\in\mptn(H)$ and $N(\alpha)$ is a clique in $H$, we have $X[N(\alpha)] = O$, which further implies $C\circ X = O$ as $H$ is a supergraph of $G - \alpha$.  
Since $B$ is zero except for the first $b$ rows, each column of $XB$ is a linear combination of the first $b$ columns of $X$.  The first $b$ columns of $X$ have their first $b$ entries zero by $X[N(\alpha)] = O$, so the first $b$ rows of $XB$ are zero.  Consequently, the first $b$ rows of $-XBQ^{-1}$ and the first $b$ columns of $-Q^{-1}B\trans X$ are zero.  Moreover, as $B\trans$ is zero except for the first $b$ columns and $XB$ is zero on the first $b$ rows, we have $B\trans XB = O$ and $Q^{-1}B\trans XBQ^{-1} = O$.  Combining all these observations, we have $A\circ\widehat{X} = I\circ\widehat{X} = O$.  

For SAP, it is straightforward to see that if $(A/Q)X = O$, then 
\[
    (E\trans AE)(O\oplus X) = (Q\oplus A/Q)(O\oplus X) = O
\]
and
\[
    A\widehat{X} = AE(O\oplus X)E\trans = (E\trans)^{-1}(E\trans AE)(O\oplus X)E\trans = O.
\]
In summary, if $A/Q\circ X = I\circ X = O$ and $(A/Q)X = O$, then $A\circ\widehat{X} = I\circ\widehat{X} = O$ and $A\widehat{X} = O$.  Suppose $A$ has SAP.  Then $\widehat{X} = O$, which implies $O\oplus X = O$ and $X = O$.  Therefore, $A/Q$ has SAP.  

For $i$-SNIP, we use the notation $\overset{(i,:]}{=} O$ to indicate that a matrix is zero except for possibly the $i$-th row.  Following a similar argument, if $(A/Q) X \overset{(i,:]}{=} O$, then $(E\trans AE)(O\oplus X) \overset{(i,:]}{=} O$, which implies $(E\trans AE)(O\oplus X)E\trans \overset{(i,:]}{=} O$.  Since each row of $(E\trans)^{-1}$ has entry $i$ equal to $0$ except possibly for the $i$-th row,
\[
    A\widehat{X} = AE(O\oplus X)E\trans = (E\trans)^{-1}(E\trans AE)(O\oplus X)E\trans \overset{(i,:]}{=} O.
\]
Therefore, if $A$ has $i$-SNIP, then $A/Q$ has $i$-SNIP.  
\end{proof}

\end{document}
